problem:  
Let $\{(M_i, g_i)\}_{i=1}^\infty$ be a sequence of compact $n$–dimensional Riemannian manifolds with  
\[
\mathrm{Ric}(M_i,g_i) \geq (n-1)k, \qquad \operatorname{diam}(M_i,g_i) \leq D
\]
for some fixed constants $k \in \mathbb{R}$ and $D>0$.  
Suppose that each $(M_i,g_i)$ admits a smooth **Riemannian foliation** $\mathcal{F}_i$ such that  
1. every leaf $L_i$ has diameter $\leq D_0$, and the second fundamental form and mean curvature of $L_i$ are uniformly bounded independent of $i$;  
2. the quotient metric space $B_i := M_i/\mathcal{F}_i$ (endowed with the leaf–space metric induced by $g_i$) is an Alexandrov space of curvature $\geq k_B$ uniformly bounded below.

Assume that $(M_i,g_i)$ converges in the Gromov–Hausdorff sense to a compact metric space $(X,d_X)$.  

**Conjecture:**  
Under the above assumptions, there exists a compact metric space $(B_\infty, d_{B_\infty})$ and a 1–Lipschitz surjection $\pi_\infty : X \to B_\infty$ with the following properties:
- $\pi_\infty$ is a **metric submetry**, i.e. $d_{B_\infty}(\pi_\infty(x),\pi_\infty(y)) = d_X(\pi_\infty^{-1}(\pi_\infty(x)), \pi_\infty^{-1}(\pi_\infty(y)))$;
- the fibers of $\pi_\infty$ arise as pointed Gromov–Hausdorff limits of the leaves of $\mathcal{F}_i$;
- $B_\infty$ is the Gromov–Hausdorff limit of the sequence $\{B_i\}$.

Equivalently, the **Riemannian foliation structure is stable under Gromov–Hausdorff convergence** for sequences satisfying uniform lower Ricci bounds and uniform control on the leaf geometry.

Why is it a "good" problem:  
This conjecture isolates a precise and unexplored aspect of geometric stability. While classical collapsing theory (Cheeger–Fukaya–Gromov) handles sectional curvature bounds, and Cheeger–Colding theory addresses Ricci-bound convergence, there is a deep gap between them: how fine geometric decompositions, like Riemannian foliations and measured fibrations, behave in purely Ricci-controlled limits.  

If true, this conjecture would generalize the Cheeger–Fukaya–Gromov fibration stability to the Ricci curvature setting, revealing new rigidity phenomena and bridging smooth and metric geometries. Its statement is minimal yet profound: it asks whether the **fibration−quotient structure itself** survives the smoothing-to-singular transition inherent to Gromov–Hausdorff limits.  

Conceptually, it unifies analytic curvature control (Ricci bounds) with geometric-topological decomposition (foliations) and metric-limit theory (Gromov–Hausdorff convergence). The question’s clarity—does the leaf/fibration structure persist in the limit?—is simply stated but demands deep new techniques, involving convergence of measured foliations, stratification of singular sets, and quantitative geometric analysis across different frameworks.